\begin{frame}{같은 데이터, 다른 갈래: ``라벨 관점''이 문제를 바꾼다}
  \small
  한 이커머스: 각 고객의 (연령, 지역, 월 구매 횟수, 평균 구매액,
  최근 90일 방문 빈도) + ``지난 30일 구독 해지 여부(1/0)'' 라벨.
  \begin{enumerate}
    \item \textbf{그냥 있는 그대로} --- $y$ = 해지 여부 예측 $\to$
      \kb{지도학습·분류}(이탈 예측 모델, Week 2$\sim$7의 도구 중 하나).
    \item \textbf{라벨만 버리면} --- 특징만으로 ``비슷한 구매 성향을
      묶어라'' $\to$ \kb{비지도학습·군집화}. 군집이 나중에 ``이 군집의
      이탈 확률이 높다''는 통찰로 이어질 수도 --- \textbf{라벨을 버린
      비지도가 라벨을 쓴 지도보다 넓은 시야}를 줄 수 있음.
    \item \textbf{한 걸음 더} --- 매 주 쿠폰(10\%/20\%/무료배송)을
      직접 결정하고, 다음 주 매출(보상)을 관찰하며 \textbf{시행착오}로
      개선 $\to$ \kb{강화학습}. ``이 고객에게 이 쿠폰이 정답''이라는
      라벨은 어디에도 없고, \textbf{보상(매출 변화)만}.
  \end{enumerate}
  \vskip0.4em
  \kq{같은 원본 데이터셋에서, 무엇을 라벨로 볼 것인가에 따라 \textbf{세
  갈래 모두가} 만들어진다.}
  \vskip0.3em
  {\footnotesize ML1의 주차가 지도(Week 2$\sim$13) + 비지도
  (Week 14$\sim$15) 두 갈래에 머무는 이유: 강화학습은 ``정답이 아예
  없고 보상만 있다'' --- 데이터의 형태 자체가 \textbf{근본적으로
  달라} 별도 과목(ML2)을 통째로 할애할 만큼 다르다.}
\end{frame}
